% Part: lambda-calculus
% Chapter: lambda-definability
% Section: primitive-recursive-functions

\documentclass[../../../include/open-logic-section]{subfiles}

\begin{document}

\olfileid{lam}{ldf}{prf}
\olsection{الدوال العودية البدائية \usetoken{S}{lambda definable}}

تذكر أن الدوال العودية البدائية هي الدوال التي يمكن تعريفها انطلاقًا
من الدوال الأساسية~$\Zero$ و$\Succ$ و$\Proj{n}{i}$ باستعمال التركيب
والعودية البدائية.

\begin{lem}
  \ollabel{lem:basic}
  الدوال العودية البدائية الأساسية~$\Zero$ و$\Succ$ والإسقاطات
  $\Proj{n}{i}$ !!{lambda definable}.
\end{lem}

\begin{proof}
هذه الدوال !!{lambda defined} بالحدود الآتية:
\begin{align*}
  \fn{Zero} & \ident \lambd[a][\lambd[fx][x]]\\
  \fn{Succ} & \ident \lambd[a][\lambd[fx][f (a f x)]] \\
  \fn{Proj}^n_i & \ident \lambd[x_0\dots x_{n-1}][x_i]
\end{align*}
\end{proof}

\begin{lem}
  \ollabel{lem:comp} افترض أن الدالة~$f$ ذات~$k$ من الوسائط، والدوال
  $g_0, \dots, g_{k-1}$ ذات~$n$ من الوسائط، !!{lambda definable}
  بالحدود~$F$ و$G_0$، \dots،~$G_{k-1}$، وأن~$h$ معرّفة منها
  بالتركيب. فعندئذ تكون~$h$ !!{lambda definable} بالحد~$H$.
\end{lem}

\begin{proof}
  يمكن أن !!{lambda define} الحد الآتي~$h$:
  \[
  H \ident \lambd[x_0 \dots x_{n-1}][F\, (G_0 x_0 \dots
    x_{n-1}) \dots (G_{k-1} x_0 \dots x_{n-1})]
  \]
  ونترك التحقق من هذه الحقيقة تمرينًا.
\end{proof}

\begin{prob}
  أكمل إثبات \olref[lam][ldf][prf]{lem:comp} ببيان أن
  $H\num{n_0}\dots\num{n_{n-1}} \red \num{h(n_0, \dots,
    n_{n-1})}$.
\end{prob}

لاحظ أن \olref{lem:comp} لم تشترط أن تكون~$f$ و$g_0$،
\dots،~$g_{k-1}$ عودية بدائية؛ وإنما اشترطت فقط أن تكون دوال كلية
!!{lambda definable}.

\begin{lem}\ollabel{lem:prim}
  افترض أن~$f$ دالة ذات~$n$ من الوسائط وأن~$g$ دالة ذات~$n+2$ من
  الوسائط، وأنهما !!{lambda definable} بالحدين~$F$ و$G$، وأن الدالة
  $h$ معرّفة من~$f$ و$g$ بالعودية البدائية. فعندئذ تكون~$h$ أيضًا
  !!{lambda definable}.
\end{lem}

\begin{proof}
  تذكر أن~$h$ معرّفة بالمعادلتين
  \begin{align*}
    h(x_1, \dots, x_n, 0) &= f(x_1, \dots, x_n)\\
    h(x_1, \dots, x_n, y+1) & = g(x_1, \dots, x_n, y, h(x_1, \dots, x_n, y)).
  \end{align*}
  وبعبارة غير صورية، يكرر التعريف بالعودية البدائية تحديث الحالة الذي
  تحدده الدالة~$g$ عدد~$y$ من المرات، بدءًا من~$f(x_1, \dots, x_n)$.
  وهذا يذكّرنا بتعريف أعداد تشيرش، فهي أيضًا معرّفة بوصفها مكررات.

  وللتبسيط نقدم التعريف والبرهان في حالة~$n=1$: وسيط واحد~$x$ يضاف
  إلى وسيط العودية~$y$. وتكون الدالة~$h$ !!{lambda defined} بالحد:
  \begin{align*}
    H \ident & \lambd[x][\lambd[y][\fn{Snd} (y
        D \tuple{\num{0}, F x})]]
    \intertext{حيث}
    D \ident & \lambd[p][\tuple{\fn{Succ} (\fn{Fst}\, p), (G x
          (\fn{Fst}\, p) (\fn{Snd}\, p))}]
  \end{align*}
  حالة التكرار التي نحفظها زوج، مركبته الأولى القيمة الحالية لـ$y$
  والثانية القيمة الموافقة لـ$h$. وفي كل خطوة تكرار ننشئ زوجًا من
  القيم الجديدة لـ$y$ و$h$، تحدده الدالة~$g$؛ وبعد انتهاء التكرار
  نعيد المركبة الثانية من الزوج، وهي القيمة النهائية لـ$h$. ونثبت
  الآن أن هذا تمثيل للعودية البدائية بالفعل.

  نريد أن نثبت، لكل~$n$ و$m$، أن~$H\,\num{n}\,\num{m} \red
  \num{h(n,m)}$. ولهذا نبين أولًا أنه إذا كان
  $D_n \ident \Subst{D}{\num{n}}{x}$، فإن
  $D_n^m \tuple{\num{0}, F\, \num{n}} \red
  \tuple{\num{m}, \num{h(n, m)}}$. ونجري الاستقراء على~$m$.

  إذا كان~$m=0$، فنريد
  $D_n^0 \tuple{\num{0}, F\, \num{n}} \red
  \tuple{\num{0}, \num{h(n, 0)}}$. لكن
  $D_n^0 \tuple{\num{0}, F\, \num{n}}$ هو نفسه
  $\tuple{\num{0}, F\, \num{n}}$. وبما أن~$F$ !!{lambda define}s~$f$،
  يختزل هذا إلى~$\tuple{\num{0}, \num{f(n)}}$، وبما أن
  $f(n) = h(n, 0)$، فهذا هو~$\tuple{\num{0}, \num{h(n,0)}}$.

  والآن افترض أن
  $D_n^m \tuple{\num{0}, F\, \num{n}} \red
  \tuple{\num{m}, \num{h(n, m)}}$. نريد أن نبين أن
  $D_n^{m+1} \tuple{\num{0}, F\, \num{n}} \red
  \tuple{\num{m+1}, \num{h(n, m+1)}}$.
  \begin{align*}
    D_n^{m+1} \tuple{\num{0}, F\, \num{n}}
    & \ident D_n(D_n^{m} \tuple{\num{0}, F\, \num{n}})\\
    & \red D_n\,\tuple{\num{m}, \num{h(n, m)}} \text{\qquad (بفرض الاستقراء)}\\
    & \ident (\lambd[p][
      \tuple{
        \fn{Succ} (\fn{Fst}\, p), (G \, \num{n} (\fn{Fst}\, p) (\fn{Snd}\, p))
    }]) \tuple{\num{m}, \num{h(n,m)}}\\
    & \redone
    \tuple{\fn{Succ} (\fn{Fst}\, \tuple{\num{m}, \num{h(n,m)}}),\\
      & \qquad (G \, \num{n} (\fn{Fst}\, \tuple{\num{m}, \num{h(n,m)}}) (\fn{Snd}\, \tuple{\num{m}, \num{h(n,m)}}))}\\
    & \red
    \tuple{\fn{Succ}\, \num{m}, (G \, \num{n} \,\num{m} \, \num{h(n,m)})}\\
    & \red \tuple{\num{m+1}, \num{g(n, m, h(n, m))}}
  \end{align*}
  وبما أن~$g(n, m, h(n, m)) = h(n, m+1)$، فقد تم البرهان.

  وأخيرًا، انظر إلى
  \begin{align*}
    H\,\num n\,\num m &\ident \lambd[x][\lambd[y][\fn{Snd} (y
        (\lambd[p].\tuple{\fn{Succ} (\fn{Fst}\, p), (G\, x\,
          (\fn{Fst}\, p)\, (\fn{Snd}\, p))}) \tuple{\num{0}, F x})]]\\
    & \qquad\,\num n\, \num m\\
    & \red \fn{Snd} (\num m \,
    \underbrace{(\lambd[p].\tuple{\fn{Succ} (\fn{Fst}\, p), (G \,\num n\,
      (\fn{Fst}\, p) (\fn{Snd}\, p))})}_{D_n} \tuple{\num{0}, F \num n})\\
    & \ident \fn{Snd} (\num m \, D_n\, \tuple{\num{0}, F \num n})\\
    & \red \fn{Snd} \,(D_n^m  \tuple{\num{0}, F \num n})\\
    & \red \fn{Snd} \,\tuple{\num{m}, \num{h(n, m)}} \\
    & \red \num{h(n,m)}.
  \end{align*}
\end{proof}


\begin{prop}
  كل دالة عودية بدائية !!{lambda definable}.
\end{prop}

\begin{proof}
  بحسب \olref{lem:basic}، جميع الدوال الأساسية !!{lambda definable}،
  وبحسب \olref{lem:comp} و\olref{lem:prim}، تنغلق الدوال
  !!{lambda definable} تحت التركيب والعودية البدائية.
\end{proof}

\end{document}
